% Status: result_ready_with_limits. Bounded-draft-ready.
\chapter{H3: Wallet-Timing}
\label{ch:h3}

H3 untersucht, ob dataset-relative Wallet-Tiers zeitliche Aktivitätsmuster vor
oder um Polymarket-Preisbewegungen zeigen. Die Tiers werden aus den beobachteten
kumulierten Betragsperzentilen der Wallet-Verteilung abgeleitet, nicht über eine
willkürliche USD-Schwelle \citep{zotero_poly_001,zotero_poly_007}. Die
Timing-Diagnostik nutzt Lead-Lag-Korrelationen und Granger-Tests
\citep{lit_granger_001}.

\paragraph{Wallet-Tiers.}
Abbildung~\ref{fig:h3tiers} zeigt die verteilungsbasierten Tiers und ihre
Besetzung. Das oberste Tier umfasst das oberste eine Prozent der Wallets
(32 Adressen im beobachteten, BUY-seitigen Datensatz). Die Tiers sind damit
relativ zum Datensatz definiert und tragen keine externe Schwelle.

\begin{figure}[ht]
\centering
\includegraphics[width=0.78\textwidth]{thesis_h3_wallet_tier_counts.png}
\caption{Dataset-relative Wallet-Tiers und ihre Besetzung. Quelle:
\texttt{data/results/thesis\_h3\_wallet\_tier\_counts.png}.}
\label{fig:h3tiers}
\end{figure}

\paragraph{Klarstes Muster im obersten Tier.}
Das oberste Tier zeigt die deutlichste deterministische Timing-Struktur: eine
Lag-1-Korrelation von $0.186$ zwischen Tier-Aktivität und nachfolgenden
Preisänderungen und einen minimalen Granger-$p$-Wert von $0.0012$ über 1216
ausgerichtete Tageszeilen (Abbildung~\ref{fig:h3granger}). Das ist ein
prädiktives Timing-Signal unter den getesteten Modellannahmen, kein Nachweis
von Kausalität oder privater Information \citep{zotero_poly_005}.

\begin{figure}[ht]
\centering
\includegraphics[width=0.78\textwidth]{thesis_h3_granger_pvalues.png}
\caption{Granger-Diagnostik je Tier (minimale $p$-Werte). Quelle:
\texttt{data/results/thesis\_h3\_granger\_pvalues.png}.}
\label{fig:h3granger}
\end{figure}

\paragraph{Deskriptives Lead-Time-Profil.}
Abbildung~\ref{fig:h3lead} stellt die Lead-Time der Tier-Aktivität relativ zu
ausgewählten Bewegungen dar. Das Profil ist rein deskriptiv und beschreibt, wie
sich beobachtete Aktivität zeitlich um Preisänderungen verteilt, ohne eine
Wirkungsrichtung zu behaupten.

\begin{figure}[ht]
\centering
\includegraphics[width=0.78\textwidth]{thesis_h3_lead_time_amount.png}
\caption{Deskriptives Lead-Time-Profil der Tier-Aktivität. Quelle:
\texttt{data/results/thesis\_h3\_lead\_time\_amount.png}.}
\label{fig:h3lead}
\end{figure}

\begin{table}[ht]
\centering
\caption{H3 --- Wallet-Tier-Timing-Diagnostik (T4)}
\label{tab:t4}
\small
\begin{tabular}{@{}lrrr@{}}
\toprule
Tier & Lag-1-Korrelation & Granger-$p$ (min) & ausgerichtete Zeilen \\
\midrule
Oberstes 1\%-Tier & $0.186$ & $0.0012$ & 1216 \\
\bottomrule
\end{tabular}
\end{table}

\section{Ereignis-Wallet-Anomalien als deskriptive Diagnostik}
\label{sec:anomalie}
Auf derselben Datenbasis prüft eine deterministische Anomalie-Diagnostik, ob sich
um die kuratierten Ereignisse auffällige Aktivitäts-, Bewegungs- oder
Konzentrationsmuster zeigen. Eine Tagesbeobachtung gilt als auffällig, wenn ihr
robuster $z$-Wert mindestens $2.0$ erreicht oder ihr Perzentilrang mindestens
$0.95$ beträgt, gemessen gegen ein Basisfenster von 30 bis 8 Tagen vor dem
Ereignis und ein Ereignisfenster von einem Tag davor bis drei Tage danach. Die
Diagnostik betrachtet vier Familien: aktive Wallets je Tier, absolute
Marktbewegung, Top-Tier-Konzentration und Tier-Betragssumme.

Als Beispiel zeigt die Verurteilung im Schweigegeld-Prozess vom 30.~Mai 2024 im
Basis-Tier eine deutliche Aktiv-Wallet-Auffälligkeit ($z = 3.5$, Perzentilrang
$1.0$, zwei auffällige Tage) sowie eine auffällige Marktbewegung
(Perzentilrang $0.96$, ein Tag). Abbildung~\ref{fig:h3anom} fasst über alle
sieben Ereignisse zusammen, welche Familien an welchen Ereignissen auffällige
Tage aufweisen.

\begin{figure}[ht]
\centering
\includegraphics[width=0.82\textwidth]{thesis_h3_event_wallet_anomalies.png}
\caption{Ereignis-Wallet-Anomalien je Ereignis und Anomalie-Familie
(deskriptiv). Quelle: \texttt{data/results/thesis\_h3\_event\_wallet\_anomalies.png}.}
\label{fig:h3anom}
\end{figure}

\paragraph{Was die Diagnostik nicht ist.}
Die Diagnostik arbeitet ausschliesslich mit aggregierten Tier-Zählungen und gibt
keine Wallet-Adressen aus. Sie ist als historische, deskriptive
Auffälligkeits-Diagnostik zu lesen und liefert Prüf-Hinweise, kein Beweis von
privater Information, Insiderhandel, Fehlverhalten, Kausalität oder
Profitabilität. Die spätere Review-Schicht, die solche Hinweise einordnet, wird
in Kapitel~\ref{ch:erweiterungen} beschrieben.

\paragraph{Limitation.}
Die Wallet-Daten enthalten nur Kauftransaktionen und sind quellseitig gefiltert,
die Ausrichtung erfolgt auf Tagesebene, und über mehrere Tiers und Lags hinweg
besteht ein Multiple-Testing-Problem. Die Granger- und Anomalie-Ausgaben sind
deshalb als prädiktive Timing- und Auffälligkeits-Diagnostik zu lesen, nicht als
Beweis von Kausalität, privater Information oder Profitabilität.

\section{Signaturen informierten Handels und ein Verdachts-Profil}
\label{sec:verdacht}
Forschung und Praxis nennen wiederkehrende Signaturen informierten oder
Insider-Handels: frische Wallets mit wenigen Lebenstransaktionen, die grosse
Positionen eingehen, ungewöhnliche Positionsgrösse vor noch nicht angekündigten
Ausgängen, abnormale Handelsgrösse und Volumen, hohe Konzentration und
einseitiger Fluss sowie Aktivität, die sich gegen die Marktauflösung verdichtet.
Akademisch zeigt \citet{zotero_delvecchio_001} für Kalshi-Märkte, dass
informierter Handel real und über eine Abnormal-Trade-Size-Kennzahl detektierbar
ist und sich gegen die Auflösung verstärkt, im Einklang mit dem Insider-Modell
von \citet{lit_kyle_001}. Praktisch setzt Polymarket eine On-Chain-Überwachung
ein, die unter anderem ungewöhnliche Positionsgrösse vor Ergebnissen und
auffällige Wallets flaggt \citep{news_polymarket_integrity_001}, gestützt auf
forensische Heuristiken wie Funding-Quellen und Wallet-Clustering
\citep{news_chainalysis_forensics_001}. Konzeptionell rahmt
\citet{zotero_poly_005} Insiderhandel in Prognosemärkten.

\paragraph{Profil aus dem getesteten Signatur-Modul.}
Dieselben Signaturen werden in einem deterministischen, getesteten Modul (\texttt{informed\_trading\_signature}) auf den BUY-seitigen Whale-Datensatz
(ab 10\,000 USD) anwenden. In den sieben Ereignisfenstern dominieren neu
erscheinende Wallets. Im Median traten 54\% der aktiven Wallets erst um das
Ereignis herum erstmals auf (Spanne 48--69\%), bei einer Top-1-Konzentration von
im Median 64\% (54--81\%) und gegenüber der Vorperiode erhöhter
Aktivität und erhöhtem Volumen (Abbildung~\ref{fig:h3profile}). Das Modul fasst die Einzel-Signaturen perzentil-normalisiert zu einem \texttt{suspicion\_score} zusammen (Spanne 0.10--0.81, Median 0.55). Dieses Muster
deckt sich mit der Literatur-Signatur aus frischen Wallets, grosser
Positionsgrösse und Konzentration.

\begin{figure}[ht]
\centering
\includegraphics[width=0.86\textwidth]{h3_informed_trading_profile.png}
\caption{Ereignisfenster-Profil aus dem getesteten Signatur-Modul: Anteil neuer Wallets,
Top-1-Konzentration und Fenster-Volumen je Ereignis (BUY-only, ab 10\,000 USD).
Quelle: \texttt{data/results/h3\_informed\_trading\_signature.csv}.}
\label{fig:h3profile}
\end{figure}

\paragraph{Verdachts-Profil als Hypothese.}
Daraus folgt ein Verdachts-Profil, das als Hypothese und Prüf-Filter dient. Ein
Ereignisfenster ist prüfwürdig, wenn mehrere Signaturen gemeinsam auftreten --
ein hoher Anteil neuer Wallets, hohe Konzentration, abnormale Aktiv-Wallet- und
Volumen-Werte aus Abschnitt~\ref{sec:anomalie} sowie ein Vorlauf der
Tier-Aktivität vor der Preisbewegung aus der Lead-Lag- und Granger-Diagnostik.
Genau diese Kombination speist die Review-Queue in
Kapitel~\ref{ch:erweiterungen}.

\paragraph{Was das Profil nicht beweist.}
Ein hoher Anteil neuer Wallets ist kein Insider-Nachweis. Aufmerksamkeitsstarke
Ereignisse ziehen allgemein neue Teilnehmer an, frische Adressen entstehen auch
aus Datenschutz- oder Komfortgründen, und der Datensatz ist BUY-only und ab
10\,000 USD gefiltert, was den Anteil neuer Wallets mechanisch erhöht. Zudem
zeigen \citet{zotero_poly_001} für die Oktober-Episode 2024, dass Kapital auf
beide Marktseiten floss, vereinbar mit heterogenen Erwartungen statt
Manipulation. Das Profil bleibt damit eine Verdachts-Hypothese und ein
Prüf-Filter, kein Nachweis privater Information.

\paragraph{Was die Diagnostik härten würde.}
Zur Härtung fehlen deterministische, getestete Merkmale: die Funding-Quelle der
Wallets (zentrale Börse oder Mixer), Wallet-Lebensdauer und Transaktionszahl,
einseitiger Fluss je Markt, der Vorlauf gegenüber verifizierten
Nachrichten-Zeitstempeln und ein Out-of-Sample-Test gegen die wenigen
dokumentierten Insider-Fälle. Diese Erweiterungen sind als Spezifikation
vermerkt und bleiben an die Audit- und Wallet-Adress-Grenzen gebunden.
